%==============================================================================
% Chapter 5 -- Future work
%==============================================================================
\chapter{Future work}
\label{ch:future-work}

% -----------------------------------------------------------------------------
% The "expanding funnel": start from a very concrete immediate next step and
% widen to a multi-year research vision, closing with a Gantt chart. Edit the
% pgfgantt block at the end to match your real milestones and dates.
% -----------------------------------------------------------------------------

Having established that whole symmetric assemblies can be modelled end to end
(Chapters~\ref{ch:methods}--\ref{ch:results}), I now describe the work that
builds on this. Section~\ref{sec:experimental-validation} concerns experimental
validation of the designs, Section~\ref{sec:motif-scaffolding} the extension
from unconditional to symmetric motif scaffolding generation, and
Section~\ref{sec:shape-control} finer control over assembly shape.
Section~\ref{sec:broad-agenda} places these in the context of the wider research
programme, and Section~\ref{sec:gantt} gives the timeline.


\section{From \emph{in silico} metrics to experimental validation}
\label{sec:experimental-validation}

The results of Chapter~\ref{ch:results} are entirely computational. The metrics
of Section~\ref{sec:evaluation} establish that the generated particles are
geometrically sound and that their designed interfaces are comparable to those
of a natural reference particle. They do not establish that a design will
express solubly, fold, or assemble into the intended cage. The usual substitute,
refolding with a structure predictor, is unreliable at this chain count
\parencite{bryant_predicting_2022}, which is why it was avoided here, and
designed assemblies that satisfy computational criteria have been found to be
structurally flexible or oligomorphic when characterised experimentally
\parencite{khmelinskaia_local_2025}. The real value of this proposed method should 
be assessed through experimental validation.

I will spend part of the 2026/27 academic year at the Baker Lab, where the relevant
experimental pipeline is well established. The work there will begin with a
small design campaign: generating a library of assemblies, designing sequences
with ProteinMPNN \parencite{dauparasRobustDeepLearning2022}, filtering on the
metrics of Section~\ref{sec:evaluation}, and expressing the candidates that
pass those filters. Characterisation would proceed in order of increasing cost (as previously done by \cite{haas_novo_2026} and coworkers), beginning
with assays like Immobilized Metal Affinity Chromatograpy (IMAC), Size Exclusion Chromatography (SEC), and chromatography and LC--MS to confirm subunit identity and an elution profile consistent with an assembled particle. These experiments will be followed by negative-stain
electron microscopy to assess morphology, and single-particle cryo-EM for the
most promising candidates. How much of this can be completed within the year
will depend on expression yields and on the number of designs that clear the
early filters.

The campaign would also yield matched computational and experimental outcomes
for assemblies generated end to end rather than docked from pre-computed
building blocks, which is not something currently available for this class of
protiens. Such data would indicate whether the inexpensive geometric filters used
in this report track successful assembly in practice, and would inform the
broader programme described in Section~\ref{sec:broad-agenda}.


\section{From unconditional to conditional generation: symmetric motif scaffolding}
\label{sec:motif-scaffolding}

Design-CP as reported performs \emph{unconditional} generation: it produces a
symmetric particle, but not a particle built around a \emph{specific} antigen. The
natural next step is \emph{conditional} generation, holding many copies of one or more functional
motifs (the antigen, or an epitope we wish to display) fixed while the model
designs a symmetric scaffold that presents them. In the symmetric setting this is
\emph{symmetric motif scaffolding}: the standard motif-scaffolding problem
\parencite{watsonNovoDesignProtein2023,didiFrameworkConditionalDiffusion2024},
recently generalised to atom-level motifs without known sequence indices
\parencite{ahern_atom_2025}, composed with the
point-group machinery of Section~\ref{sec:rfd3-prelim}, so that a motif placed in
one ASU (or at symmetric interface between some ASUs) is presented at every symmetry-equivalent 
position on the particle. This work will be
carried out in collaboration with the Baker Lab during a visit there in the next
academic year (2026/27), giving direct access to their design and experimental
pipelines for validating conditionally designed immunogens.

\section{Greater control over the shape of generated assemblies}
\label{sec:shape-control}


Particle radius is already controllable within the framework of
Chapter~\ref{ch:methods}, but there are other geometric features that a protein designer might want to design for. Two are of
particular interest. First, the currently generated assemblies show pores at some interfaces that arise as a by-product of the design rather than as a specified feature. Second, natural capsid subunits frequently possess extended arms that interdigitate with neighbouring subunits, a feature with clear structural
consequences: the N-terminal arm forms the $\beta$-annulus underlying
quasi-equivalence in $T{=}3$ plant viruses \parencite{harrison_protein_2017}, and
C-terminal arm exchange allows papillomavirus and polyomavirus pentamers to
occupy hexavalent positions in a $T{=}7$ lattice
\parencite{liddingtonStructureSimianVirus1991}. Since the training distribution is dominated
by globular monomers, such subunits fall outside what the model readily
produces.

Both features are geometric, which suggests shape conditioning as a route to
controlling them. Building on existing spatial conditioning methods
\parencite{ingrahamIlluminatingProteinSpace2023,stark_protcomposer_2025,faltings_proxelgen_2025,raghu_multiscale_2025,li_robust_2026},
I plan to let the designer specify a target envelope that the model must
respect, which would make pore size and position explicit design variables and
allow non-globular subunits to be requested directly.


\section{A broader agenda: designing large symmetric protein assemblies}
\label{sec:broad-agenda}

More broadly, and over the remainder of the DPhil, I will continue to study the
design of large, and often symmetric, protein assemblies as a subject in its own
right. Doing this well requires understanding, and learning to design for, the
factors that govern whether subunits actually come together: specific
protein--protein interactions at the designed interfaces and the
thermodynamics of self-assembly, the functional and physicochemical drivers of
symmetric oligomerisation \parencite{goodsell_structural_2000} and the multistep,
ordered assembly pathways documented for natural complexes
\parencite{marsh_structure_2015}, and phenomena such as the structural
flexibility and oligomorphism observed in designed assemblies
\parencite{khmelinskaia_local_2025}.
The space of targets is correspondingly wide, from vaccine scaffolds to
mechanically coupled molecular machines \parencite{courbet_computational_2022} and
allosterically switchable assemblies \parencite{pillai_novo_2024}.

\section{Timeline and milestones}
\label{sec:gantt}

Figure~\ref{fig:gantt} summarises the plan on a single timeline: the work
packages of Sections~\ref{sec:motif-scaffolding}--\ref{sec:broad-agenda}, the
papers and conferences they are aimed at, and the two formal milestones of the
DPhil.

\clearpage % put the Gantt chart on its own page so it can't overlap the text above
\begin{figure}[H]
  \centering
  \hspace*{-1.2cm}% shift the chart left (tweak as needed)
  \begin{ganttchart}[
      hgrid, vgrid,
      x unit=1.45cm,
      y unit title=0.6cm,
      y unit chart=0.75cm,
      title height=1,
      title/.append style={fill=oxfordblue!12},
      title label font=\footnotesize\bfseries,
      bar/.append style={fill=oxfordblue!75, rounded corners=1.5pt},
      bar height=0.55,
      bar label font=\footnotesize,
      bar label node/.append style={text width=5.2cm, align=right},
      group/.append style={fill=oxfordblue},
      group height=0.55,
      group label font=\footnotesize\bfseries,
      group label node/.append style={text width=5.2cm, align=right},
      milestone/.append style={fill=orange!85},
      milestone height=0.55,
      milestone label font=\footnotesize\itshape,
      milestone label node/.append style={text width=5.2cm, align=right},
      group right shift=0,
      group left shift=0,
    ]{1}{8}
    % --- title rows: years and quarters (TODO: adjust to your real timeline) ---
    \gantttitle{2026}{2}\gantttitle{2027}{4}\gantttitle{2028}{2} \\
    \gantttitle{Q3}{1}\gantttitle{Q4}{1}%
    \gantttitle{Q1}{1}\gantttitle{Q2}{1}\gantttitle{Q3}{1}\gantttitle{Q4}{1}%
    \gantttitle{Q1}{1}\gantttitle{Q2}{1} \\
    % --- Foundation: whole-assembly modelling ---
    \ganttgroup{Whole-assembly modelling}{1}{5} \\
    \ganttbar{Finalise Design-CP \& submit paper}{1}{2} \\
    \ganttbar{Experimental (in-vitro) validation}{3}{5} \\
    % --- Conditional generation ---
    \ganttgroup{Conditional generation}{3}{8} \\
    \ganttbar{Symmetric motif scaffolding}{3}{6} \\
    \ganttbar{Test on real antigens}{6}{8} \\
    % --- Shape / spatial control ---
    \ganttgroup{Shape \& spatial control}{1}{4} \\
    \ganttbar{Improve spatial/shape\\conditioning}{1}{4} \\
    % --- Broader agenda ---
    \ganttgroup{Large symmetric assemblies}{6}{8} \\
    \ganttbar{PPIs \& thermodynamics\\of assembly}{6}{8} \\
    \ganttbar{Thesis writing}{7}{8} \\
    % --- Papers and conferences ---
    % Each paper bar ends on the submission deadline of the conference it targets:
    % ICML 2027 closes in 2027 Q1, ICLR 2028 in 2027 Q3.
    \ganttgroup{Papers \& conferences}{2}{8} \\
    \ganttbar{Shape-conditioning paper\\(ICML 2027)}{2}{3} \\
    \ganttbar{Symmetric-scaffolding paper\\(ICLR 2028)}{4}{5} \\
    \ganttmilestone{ICML 2027}{4} \\
    \ganttmilestone{ICLR 2028}{7} \\
    % --- Milestones ---
    \ganttmilestone{Transfer of Status}{1} \\
    \ganttmilestone{Confirmation of Status}{5}
  \end{ganttchart}
  \caption{Gantt chart of upcoming milestones and deadlines}
  \label{fig:gantt}
\end{figure}
